\documentclass[10pt,a4paper]{amsart}
\usepackage[margin=19mm]{geometry}
\usepackage{amsmath,amssymb,mathtools}
\usepackage[scr=rsfs,cal=euler]{mathalfa}
\usepackage{xfrac}
\usepackage[unicode,hidelinks,bookmarks=false]{hyperref}
\newcommand{\ie}{i.e.\ }
\newcommand{\cf}{cf.\ }
\newcommand{\resp}{respectively\ }
\newcommand{\Subsection}{\subsection}
\providecommand{\nobreakdash}{\nobreak\hbox{-}}
\setlength{\emergencystretch}{2em}
\multlinegap0pt
\usepackage{amssymb}
\usepackage{float}
\usepackage{tikz}
\newtheorem{theorem}{Theorem}
\newtheorem{corollary}{Corollary}
\title{Spectral Properties of Zero-Divisor Graphs of Truncated Polynomial Rings}
\author{Bilal Ahmad Rather}
\date{Source: arXiv:2604.03101v1 (CC BY 4.0)}
\begin{document}
\maketitle

\noindent\emph{Source and licence.} Spectral Properties of Zero-Divisor Graphs of Truncated Polynomial Rings. Version arXiv:2604.03101v1 (CC BY 4.0), distributed by arXiv under the Creative Commons Attribution 4.0 International licence, http://creativecommons.org/licenses/by/4.0/, as recorded in the arXiv licence metadata for this exact version and shown on the abstract page https://arxiv.org/abs/2604.03101v1. Full licence notice: \texttt{licenses/arxiv-2604.03101-CC-BY-4.0.txt}.\par\smallskip

\noindent\textit{Editorial scope.} Counted unit: the article's corollary even (source0.tex lines 806--814). The article's complete original proof of it is delivered with it (source0.tex lines 815--1122, one \texttt{proof} environment of 308 lines). No mathematics is written, repaired or re-derived: the statement and the proof are the article's own lines, in the article's own notation, and no retained line is substituted or re-flowed.  Retained uncounted as the source-local context the proof needs: lines 455--500.  Nothing was omitted from any retained span, so the package carries no omission record.  No retained line contains a citation, so the wrapper carries no bibliography and no bibliography entry.  No retained line contains a figure environment, an image-inclusion command or drawing code, so no image is delivered and the package declares no asset dependency.  Editorial formatting only, in addition: an AMS \texttt{amsart} preamble that restates the article's own declarations and macros used by the retained lines, namely the environments \texttt{theorem}, \texttt{corollary} on separate counters, together with the standard packages those lines need.  The retained environments print in this document as Theorem 1, Corollary 1 in order of appearance, whereas the article prints the same items with the numbering of its own full document.  The complete native source of the article as distributed in the arXiv e-print for this exact version is bundled unchanged as \texttt{sources/197787/source0.tex}.
\begin{theorem}\label{spectra thm even}
	Let $p$ be a prime and $a=2b$ with $b\in \mathbb{N}$, and let $ \varGamma(R)$ be the zero-divisor graph of $R = \mathbb{Z}_p[x]/\langle x^{2b}\rangle \cong \mathbb{F}_p[x]/(x^{2b})$. 	Then the $A_\alpha$--spectrum of $ \varGamma(R)$ is as follows: For each $i=1,\dots,2b-1$, there is an eigenvalue
		\[
		\lambda_i =
		\begin{cases}
			\alpha\bigl(p^{ i}-1\bigr), & \text{if } 1\le i \le b-1,\\[3pt]
			\alpha\bigl(p^{ i}-1\bigr) - 1, & \text{if } b\le i \le 2b-1,
		\end{cases}
		\]
		with multiplicity $n_i - 1 = (p-1)p^{ 2b-1-i} - 1$. The remaining $2b-1$ eigenvalues are the eigenvalues
		$\theta_1,\dots,\theta_{2b-1}$ of the $(2b-1)\times(2b-1)$ matrix
		$B=(B_{ij})_{1\le i,j\le 2b-1}$, where
		\[
		B = \alpha D^\ast + (1-\alpha)Q,\quad \text{with}\quad
		D^\ast = \operatorname{diag}(d_1,\dots,d_{2b-1})
		\]
		and $Q=(Q_{ij})$ the quotient matrix of $A( \varGamma(R))$
		with respect to the partition $V( \varGamma(R))=\bigcup_{i=1}^{2b-1}V_i$, given by
		\[
		Q_{ij} =
		\begin{cases}
			0, & \text{if } i=j,\ 1\le i\le b-1\ (2i<2b),\\[3pt]
			n_i - 1, & \text{if } i=j,\ b\le i\le 2b-1\ (2i\ge 2b),\\[3pt]
			n_j, & \text{if } i\ne j,\ i+j\ge 2b,\\[3pt]
			0, & \text{if } i\ne j,\ i+j<2b.
		\end{cases}
		\]
		Equivalently,
		\[
		B_{ij} =
		\begin{cases}
			\alpha d_i, & \text{if } i=j,\ 1\le i\le b-1,\\[3pt]
			\alpha d_i + (1-\alpha)(n_i-1), & \text{if } i=j,\ b\le i\le 2b-1,\\[3pt]
			(1-\alpha)n_j, & \text{if } i\ne j,\ i+j\ge 2b,\\[3pt]
			0, & \text{if } i\ne j,\ i+j<2b,
		\end{cases}
		\]
		where for $1\le i\le 2b-1$, the degree of any $v\in V_i$ is
		\[
		d_i = \deg(v) =
		\begin{cases}
			p^{ i} - 1, & \text{if } 1\le i \le b-1 \ (2i<2b),\\[3pt]
			p^{ i} - 2, & \text{if } b\le i \le 2b-1 \ (2i\ge 2b).
		\end{cases}
		\]
\end{theorem}

\begin{corollary}\label{L integral even}
	The Laplacian spectrum of $ \varGamma(R)$ is integral, and its spectrum is 
	\[ \{0\}
	\cup
	\bigcup_{\substack{1\le i\le 2b-1\\ i\neq b}}
	\bigl\{(p^{ i}-1)^{[n_i]}\bigr\}
	\cup
	\bigl\{(p^{ b}-1)^{[n_b-1]}\bigr\}. \]
\end{corollary}

\begin{proof}
	Since, for any $\alpha\neq\beta$,
	\[
	\frac{1}{\alpha-\beta}\bigl(A_\alpha(G)-A_\beta(G)\bigr)
	= D(G)-A(G)=L(G).
	\]
	So, by Theorem \ref{spectra thm even}, we get $$\alpha\bigl(p^{ i}-1\bigr)-\beta\bigl(p^{i}-1\bigr)=(\alpha-\beta)(p^{i}-1).$$
	Thus, for $1\le i \le b-1,$ it follows that $p^{i}-1$  is the Laplacian eigenvalue of $ \varGamma(R)$ with $n_i - 1 = (p-1)p^{ 2b-1-i} - 1$. Similarly, for $b\le i \le 2b-1,$ we see that 
	$$\alpha\bigl(p^{ i}-1\bigr) - 1-\beta\bigl(p^{ i}-1\bigr) + 1=(\alpha-\beta)(p^{i}-1) $$ is the eigenvalue of $ \varGamma(R)$ with multiplicity $n_i - 1 = (p-1)p^{ 2b-1-i} - 1$.  By Theorem~\ref{spectra thm even}, the quotient of $A_\alpha( \varGamma(R))$ 
	with respect to the partition $\{V_i\}_{i=1}^{2b-1}$ is
	\[
	B(\alpha) = \alpha D^\ast + (1-\alpha)Q^\ast,
	\]
	where $D^\ast=\operatorname{diag}(d_1,\dots,d_{2b-1})$ and $Q^\ast$ is the quotient of $A(G)$. Hence for any $\alpha\neq\beta$,
	\[
	\frac{1}{\alpha-\beta}\bigl(B(\alpha)-B(\beta)\bigr)
	= D^\ast - Q^\ast =: \overline{L}.
	\]
	The $(2b-1)\times(2b-1)$ matrix $\overline{L}$ is therefore the quotient of the Laplacian $L(G)$ with respect to
	the partition $\{V_i\}$. With the description of $Q^\ast$ and $d_i$ from Theorem~\ref{spectra thm even}, we have
	\[
	\overline{L}_{ij} =
	\begin{cases}
		d_i, & i=j,\ 1\le i\le b-1,\\
		d_i-(n_i-1)=p^{ i}-1-n_i, & i=j,\ b\le i\le 2b-1,\\
		- n_j, & i\ne j,\ i+j\ge 2b,\\
		0, & i\ne j,\ i+j<2b.
	\end{cases}
	\]
In explicit matrix form, $\overline{L}$ is the $(2b-1)\times(2b-1)$ matrix
\[
\overline{L}=
\begin{pmatrix}
	d_1 & 0   & 0   & \cdots & 0         & 0         & -n_{2b-1} \\[3pt]
	0   & d_2 & 0   & \cdots & 0         & -n_{2b-2} & -n_{2b-1} \\[3pt]
	0   & 0   & d_3 & \cdots & -n_{2b-3} & -n_{2b-2} & -n_{2b-1} \\[3pt]
	\vdots & \vdots & \vdots & \ddots & \vdots    & \vdots    & \vdots \\[3pt]
	0   & 0   & -n_3 & \cdots &
	p^{ 2b-3}-1-n_{2b-3} & -n_{2b-2} & -n_{2b-1} \\[3pt]
	0   & -n_2 & -n_3 & \cdots & -n_{2b-3} &
	p^{ 2b-2}-1-n_{2b-2} & -n_{2b-1} \\[3pt]
	-n_1 & -n_2 & -n_3 & \cdots & -n_{2b-3} & -n_{2b-2} &
	p^{ 2b-1}-1-n_{2b-1}
\end{pmatrix}.
\]
Now, we show that see that $\{0\}$ along with $p^{i}-1$ are the eigenvalue of $\overline{L}$ for $1\leq i\leq 2b-1$ except $i=b.$

	We let $N=2b-1$ and recall that $n_j=(p-1)p^{N-j}$. So, with 	$1\le a\le b\le N$, w ehave
	\begin{align}
		\sum_{j=a}^{b} n_j
		&= (p-1)\sum_{j=a}^{b} p^{N-j}
		= (p-1)\sum_{t=N-b}^{N-a} p^t \notag\\
		&= (p-1) \frac{p^{N-a+1}-p^{N-b}}{p-1}
		= p^{N-a+1}-p^{N-b}. \label{eq:general-sum}
	\end{align}
	In particular, for $r\ge1$, we have
	\begin{equation}\label{eq:tail-sum}
		\sum_{j=N-r+1}^{N} n_j = p^r-1.
	\end{equation}
	Due to positive semidefinite property of the Laplacian of $ \varGamma(R)$, $0$ is the simple eigenvalue of $\overline{L}$. Next, for eigenvalues $p^k-1$ with $1\le k\le b-1$, we fix an integer $k$,  $1\le k\le b-1$ and let $
	\lambda_k = p^k-1.$ Now, consider a vector  $v^{(k)}\in\mathbb{R}^N$ with 
	\begin{equation}\label{eq:v-k}
		v^{(k)}_j =
		\begin{cases}
			0, & 1\le j<k,\\
			- \varGamma_k, & j=k,\\
			1, & k<j\le N-k,\\
			0, & N-k<j\le N,
		\end{cases}
	\end{equation}
	where $ \varGamma_k$ is a scaler, we are to be chosen so that eigenequation is satisfied. For the computation of $(\overline{L}v^{(k)})_i$, we consider the following cases. For  $1\le i<k$, we have $v_i^{(k)}=0$. If $j$ is such that $v_j^{(k)}\ne0$, then 
	$k\le j\le N-k$, so it gives
	\[
	i+j\le (k-1)+(N-k) = N-1 < N+1.
	\]
	Hence $i+j\ge N+1$ never holds, so all off-diagonal entries in row $i$ of $\overline{L}$
	multiply zeros of $v^{(k)}$. Thus, we have
	\[
	(\overline{L}v^{(k)})_i = d_i v^{(k)}_i = 0 = \lambda_k v^{(k)}_i.
	\]
	For $i=k$, we have $v^{(k)}_k=- \varGamma_k$. For $j$ in the support of $v^{(k)}$ with
	$j\ne k$, we have $k<j\le N-k$, so it gives
	\[
	k+j \le k+(N-k) = N < N+1.
	\]
	Thus $k+j\ge N+1$ never holds, and there are no nonzero off-diagonal
	contributions in row $k$. So, $(\overline{L}v^{(k)})_k = d_k v^{(k)}_k.$
	Since $1\le k\le b-1$, we have $d_k=p^k-1=\lambda_k$, and we obtain 
	\[
	(\overline{L}v^{(k)})_k = (p^k-1)v^{(k)}_k
	= \lambda_k v^{(k)}_k,
	\]
	which is independent of $ \varGamma_k$. If $k<i\le N-k$, then in this range $v^{(k)}_i=1$. The only indices $j$ with $v_j^{(k)}\ne 0$	are $j=k$ and $k<j\le N-k$. However, the column $k$ does not appear in the off-diagonal part of row $i$. 	Indeed, for $j=k$, we obtain
	\[
	i+k \le (N-k)+k = N < N+1.
	\]
	So, $i+k\ge N+1$ fails and it gives $(\overline{L})_{ik}=0$. Thus $ \varGamma_k$ plays
	no role in this range. The nonzero off-diagonal terms come from those
	$j$ with $k<j\le N-k, $ with $ i+j\ge N+1.$
	The inequality $i+j\ge N+1$ is $j\ge N+1-i$, so the relevant $j$ are
	\[
	j\in\bigl[\max(k+1,N+1-i), N-k\bigr].
	\]
	Since $i\le N-k$, we get $N+1-i \ge N+1-(N-k)=k+1,$ 	so the lower bound is $j\ge N+1-i$. In summary, the contributing
	indices are
	\[
	j = N+1-i,N+2-i,\dots,N-k.
	\]
	These all satisfy $j>k$, hence $v^{(k)}_j=1$. 	If $i<b$ then $N+1-i > i$, so $j=i$ is not in this interval and the
	sum includes no diagonal index. If $i\ge b$ then $2i\ge 2b=N+1$, so
	$N+1-i \le i$, and $j=i$ \emph{is} in the interval. In this case we
	must omit $j=i$ from the off-diagonal sum. Both these cases can be treated
	uniformly by writing
	\[
	(\overline{L}v^{(k)})_i
	= d_i - \sum_{j=N+1-i}^{N-k} n_j + \varepsilon_i n_i,
	\]
	where $\varepsilon_i=0$ if $i<b$ and $\varepsilon_i=1$ if $i\ge b$. By \eqref{eq:general-sum} (with $a=N+1-i$, $b=N-k$),
	\[
	\sum_{j=N+1-i}^{N-k} n_j
	= p^{N-(N+1-i)+1}-p^{N-(N-k)}
	= p^i-p^k.
	\]
	
	If $k<i\le b-1$, then $d_i=p^i-1$ and $\varepsilon_i=0$.
	Thus, we have
	\[
	(\overline{L}v^{(k)})_i
	= (p^i-1) - (p^i-p^k)
	= p^k-1
	= \lambda_k v^{(k)}_i.
	\]
	If $b\le i\le N-k$, then $d_i=p^i-1-n_i$ and
	$\varepsilon_i=1$, so we obtain 
	\[
	(\overline{L}v^{(k)})_i
	= (p^i-1-n_i) - (p^i-p^k) + n_i
	= p^k-1
	= \lambda_k v^{(k)}_i.
	\]
	Thus the eigenvalue equation holds for all $k<i\le N-k$ for every
	choice of $ \varGamma_k$. With $N-k < i\le N$, here $v^{(k)}_i=0$, so we need $(\overline{L}v^{(k)})_i=0$.
	Now $v^{(k)}_j\ne 0$ only for $k\le j\le N-k$. Since
	$i\ge N-k+1$, for any such $j$ we have
	\[
	i+j \ge (N-k+1)+k = N+1.
	\]
	So, all of these $j$ appear in the off-diagonal sum. Therefore, we have
	\[
	(\overline{L}v^{(k)})_i
	= -\sum_{j=k}^{N-k} n_j v^{(k)}_j
	= -\Bigl(n_k(- \varGamma_k) + \sum_{j=k+1}^{N-k} n_j\cdot1\Bigr)
	=  \varGamma_k n_k - \sum_{j=k+1}^{N-k} n_j.
	\]
	This expression is independent of $i$ in this range. Thus all the
	bottom rows impose the same scalar condition $	 \varGamma_k n_k = \sum_{j=k+1}^{N-k} n_j.$
	In order to  solve  $ \varGamma_k$ explicitly , we use \eqref{eq:general-sum},
	\[
	\sum_{j=k+1}^{N-k} n_j
	= \sum_{j=k}^{N-k} n_j - n_k
	= \bigl(p^{N-k+1}-p^{N-(N-k)}\bigr) - n_k
	= (p^{N-k+1}-p^k) - n_k.
	\]
	But $n_k=(p-1)p^{N-k}$, so $	p^{N-k+1}-p^k-n_k = p^{N-k}-p^k.$
	Hence, with these facts, we obtain $	\sum_{j=k+1}^{N-k} n_j = p^{N-k}-p^k.$ Also, $n_k = (p-1)p^{N-k}$, so we have
	\[
	 \varGamma_k
	= \frac{\sum_{j=k+1}^{N-k} n_j}{n_k}
	= \frac{p^{N-k}-p^k}{(p-1)p^{N-k}}
	= \frac{1-p^{2k-N}}{p-1}
	= \frac{1-p^{ 2k-2b+1}}{p-1}.
	\]
	With the above value of $ \varGamma_k$, all bottom rows give $(\overline{L}v^{(k)})_i=0=\lambda_k v^{(k)}_i$. Therefore, we have 
	\[
	\overline{L}v^{(k)} = \lambda_k v^{(k)},\qquad
	\lambda_k=p^k-1,\quad 1\le k\le b-1.
	\]
	
	\medskip
	
	Now, for the eigenvalues $p^k-1$ of $\overline{L}$ for $b+1\le k\le N$, we fix an integer $k$ with $b+1\le k\le N$ and set again	$\lambda_k=p^k-1$. For $m = N+1-k = 2b-k,$ with $1\le m\le b-1<k$, we define $w^{(k)}\in\mathbb{R}^N$ by
	\begin{equation}\label{eq:w-k}
		w^{(k)}_j =
		\begin{cases}
			0,          & 1\le j<m,\\
			-\delta_k,  & m\le j\le k-1,\\
			1,          & j=k,\\
			0,          & k<j\le N,
		\end{cases}
	\end{equation}
	where $\delta_k$ is the  scalar, which we are required to find that $w^{(k)}$ is eigenvector of some eigenvalue of $\overline{L}$. Again we compute $(\overline{L}w^{(k)})_i$ by different ranges over $i$. For $1\le i<m$, we have $w^{(k)}_i=0$. For any $j$ with $w^{(k)}_j\ne0$, we have$m\le j\le k$, and thus
	\[
	i+j \le (m-1)+k = (N+1-k-1)+k = N.
	\]
	Hence $i+j\ge N+1$ never holds, so the off-diagonal sum of $\overline{L}$ is empty,  and we obtain
	\[
	(\overline{L}w^{(k)})_i = d_i w^{(k)}_i = 0 = \lambda_k w^{(k)}_i.
	\]
	For $i=k$, we have $w^{(k)}_k=1$. The only nonzero $w^{(k)}_j$ with $j\ne k$ occur
	for $m\le j\le k-1$, and for all such $j$,  we have
	\[
	k+j \ge k+m = k+(N+1-k)=N+1,
	\]
	and they all contribute to the off-diagonal sum. Thus, we obtain
	\[
	(\overline{L}w^{(k)})_k
	= d_k\cdot 1 - \sum_{j=m}^{k-1} n_j(-\delta_k)
	= d_k + \delta_k\sum_{j=m}^{k-1} n_j.
	\]
	As $k\ge b+1$, so $k\ge b$ and hence $d_k=p^k-1-n_k$. In order for $(\overline{L}w^{(k)})_k=\lambda_k w^{(k)}_k=p^k-1$ to hold, so we must have
	\[
	p^k-1-n_k + \delta_k\sum_{j=m}^{k-1}n_j = p^k-1,
	\]
	which is equivalent to $\delta_k\sum_{j=m}^{k-1}n_j = n_k.$
	Using \eqref{eq:general-sum} with $a=m$, $b=k-1$, we obtain
	\[
	\sum_{j=m}^{k-1} n_j
	= p^{N-m+1}-p^{N-(k-1)}
	= p^{k}-p^{N-k+1}.
	\]
	Thus
	\[
	\delta_k
	= \frac{n_k}{p^k-p^{N-k+1}}
	= \frac{(p-1)p^{N-k}}{p^k-p^{N-k+1}}= \frac{p-1}{p^{2k-N}-p}
	= \frac{p-1}{p\bigl(p^{2(k-b)}-1\bigr)}.
	\]
	With the choice of $\delta_k$, for $m\le i\le k-1$, we have  $w^{(k)}_i=-\delta_k$. The nonzero $w^{(k)}_j$ occur for $j\in[m,k-1]$ (where $w_j=-\delta_k$) and $j=k$ (where $w_k=1$). With	$m\le i\le k-1$, we have $i+k \ge m+k = N+1,$	so $j=k$ always appears in the off-diagonal sum. Among $m\le j\le k-1$, the condition $i+j\ge N+1$ is $j\ge N+1-i$. Since
	\[
	N+1-i \ge N+1-(k-1) = N-k+2 = m+1,
	\]
	we have $N+1-i \ge m$ and the contributing $j$ satisfy $j\in [N+1-i,k-1].$ 	We must again check whether $i$ lies in this interval. If $m\le i<b$, then  $N+1-i > i$, and it shows that  $i$ is not in $[N+1-i,k-1]$. In this case
	\[
	(\overline{L}w^{(k)})_i
	= d_i(-\delta_k) - \sum_{j=N+1-i}^{k-1} n_j(-\delta_k) - n_k\cdot 1
	= -\delta_k d_i + \delta_k\sum_{j=N+1-i}^{k-1} n_j - n_k.
	\]
	Using \eqref{eq:general-sum} with $a=N+1-i$, $b=k-1$, we have
	\[
	\sum_{j=N+1-i}^{k-1} n_j
	= p^{N-(N+1-i)+1}-p^{N-(k-1)}
	= p^{i}-p^{N-k+1}.
	\]
	Since $i<b$, $d_i=p^i-1$, so
	\begin{align*}
		(\overline{L}w^{(k)})_i
		&= -\delta_k(p^i-1) + \delta_k(p^i-p^{N-k+1}) - n_k= -\delta_k(p^{N-k+1}-1) - n_k.
	\end{align*}
	On the other hand, we have
	\[
	\lambda_k w^{(k)}_i = (p^k-1)(-\delta_k) = -\delta_k(p^k-1).
	\]
	Thus the eigenvalue equation is equivalent to
	\[
	-\delta_k(p^{N-k+1}-1) - n_k = -\delta_k(p^k-1)
	\]
	or
	\[
	\delta_k\bigl(p^k-p^{N-k+1}\bigr) = n_k.
	\]
	But this is precisely how $\delta_k$ was chosen above. Hence
	$(\overline{L}w^{(k)})_i=\lambda_k w^{(k)}_i$ for $m\le i<b$. Now, for $b\le i\le k-1$, we get $2i\ge 2b=N+1$. So, $N+1-i\le i$ and the interval
	$[N+1-i,k-1]$ contains $i$. Thus, we have
	\[
	\sum_{\substack{m\le j\le k-1\\ i+j\ge N+1}} n_j(-\delta_k)
	= -\delta_k\left(\sum_{j=N+1-i}^{k-1} n_j - n_i\right),
	\]
	and
	\[
	(\overline{L}w^{(k)})_i
	= d_i(-\delta_k) -\Bigl(-\delta_k\bigl(\sum_{j=N+1-i}^{k-1} n_j - n_i\bigr)\Bigr)
	- n_k
	= -\delta_k d_i + \delta_k\sum_{j=N+1-i}^{k-1}n_j - \delta_k n_i - n_k.
	\]
	Again with $\sum_{j=N+1-i}^{k-1}n_j=p^i-p^{N-k+1}$ and
	$d_i=p^i-1-n_i$, we obtain
	\begin{align*}
		(\overline{L}w^{(k)})_i
		&= -\delta_k(p^i-1-n_i) + \delta_k(p^i-p^{N-k+1}) - \delta_k n_i - n_k= \delta_k - \delta_k p^{N-k+1} - n_k\\
		&= -\delta_k(p^{N-k+1}-1) - n_k.
	\end{align*}
	As in the previous case this equals $-\delta_k(p^k-1)$ exactly when
	$\delta_k(p^k-p^{N-k+1})=n_k$, that is, for our $\delta_k$, and hence for $b\le i\le k-1$, we obtain
	\[
	(\overline{L}w^{(k)})_i = \lambda_k w^{(k)}_i.
	\]
	For the last range $i>k$, we have $w^{(k)}_i=0$. For such $i$ and with any $j\in[m,k]$, we have
	\[
	i+j \ge (k+1)+m = (k+1)+(N+1-k)=N+2>N+1.
	\]
	So, all indices in the support $\{m,\dots,k\}$ contribute to the off-diagonal sum. Thus, with the defining relation for $\delta_k$, we obtain
	\begin{align*}
		(\overline{L}w^{(k)})_i
		&= -\sum_{j=m}^{k-1} n_j(-\delta_k) - n_k\cdot1\\
		&= \delta_k\sum_{j=m}^{k-1}n_j - n_k
		= \delta_k(p^k-p^{N-k+1}) - n_k
		= 0.
	\end{align*}
	Since $\lambda_k w^{(k)}_i=0$
	as well, the eigenvalue equation holds. Thus,  it shows that for $b+1\le k\le N$,
	\[
	\overline{L}w^{(k)} = (p^k-1)w^{(k)}.
	\]
	The numbers $p^k-1$ are strictly increasing in $k$, so the
	$2b-2$ values $\{p^k-1:1\le k\le N,\ k\ne b\}$ are all distinct, and
	together with $0$ they give $2b-1=N$ distinct eigenvalues. Thus, for every $i$ with $1\le i\le 2b-1$ and $i\neq b$,
	$p^i-1$ is an eigenvalue of $\overline{L}$.
\end{proof}

\end{document}
